\documentclass{../template/myrpt}
\ExamName{数字示波器的调整和应用}

\newcolumntype{Y}{>{\centering\arraybackslash}X}
\newcolumntype{C}[1]{>{\centering\arraybackslash}m{#1}}
\newcommand{\imagecell}[1]{\parbox[c][#1][c]{\linewidth}{\centering\strut}}
\newcommand{\imagehead}[2]{\parbox[c][#1][c]{\linewidth}{\centering #2}}
\pgfplotsset{compat=1.18}

\begin{document}
\maketitle

\section{实验目的}
\begin{enumerate}[label=\arabic*. ]
  \item 了解数字示波器的基本结构、工作原理和常用功能；
  \item 掌握数字示波器显示波形、光标测量和菜单自动测量的方法；
  \item 学习使用函数发生器输出不同频率、幅值和形状的周期信号；
  \item 学会利用李萨如图形测量未知信号频率；
  \item 观察数字示波器数学运算功能对双通道信号的处理效果。
\end{enumerate}

\section{实验仪器}

本实验使用的主要仪器包括数字示波器、示波器检测线和函数发生器。

\section{实验原理}

\subsection{示波器及其基本功能}

示波器是一种把随时间变化的电压信号转换为可视波形的电子测量仪器。通常情况下，示波器屏幕的横轴表示时间，纵轴表示电压，因此可直接观察周期信号的波形、幅值、周期、频率、相位和稳定性等特征。通过调节垂直灵敏度、时基、触发电平和耦合方式，可以使被测信号稳定地显示在屏幕上，并便于进行定量测量。

按照信号处理方式，示波器可分为模拟示波器和数字示波器。模拟示波器主要利用电子束偏转直接显示输入信号；数字示波器则先对输入模拟信号进行采样和模数转换，再经过存储、运算和显示处理得到波形。数字示波器通常具有自动设置、光标测量、自动参数测量、波形存储、数学运算和数据导出等功能，适合对周期信号和瞬态信号进行较方便的测量与分析。

\subsection{数字示波器的测量方法}

数字示波器测量周期信号时，常用光标测量和菜单自动测量两种方式。光标测量由操作者移动时间光标或电压光标，读出两光标间的时间差 $\Delta t$ 或电压差 $\Delta U$，再得到周期、频率或峰峰值等参数。自动测量则由示波器根据采样数据和内部算法直接给出峰峰值、周期、频率、有效值等测量结果。

对周期信号，频率与周期满足
\[
  f = \frac{1}{T}.
\]
若信号的最大值和最小值分别为 $U_{\max}$ 和 $U_{\min}$，则峰峰值为
\[
  U_{\mathrm{pp}} = U_{\max} - U_{\min}.
\]
光标测量直观、便于理解，但会受到人工定位和读数的影响；自动测量操作快速、重复性较好，但结果依赖示波器采样、触发和内部识别算法。

\subsection{李萨如图形测频率}

当数字示波器工作在 $X$-$Y$ 显示方式时，若将一路正弦信号接入 $X$ 方向、另一路正弦信号接入 $Y$ 方向，屏幕上会显示由两路信号合成的李萨如图形。若两路信号频率成简单整数比，图形将呈稳定闭合曲线；若频率比不是简单整数比，图形会不断变化而难以稳定。

设 $f_x$ 为接入 $X$ 方向的信号频率，$f_y$ 为接入 $Y$ 方向的信号频率。若稳定李萨如图形与水平外切线的切点数为 $n_x$，与竖直外切线的切点数为 $n_y$，则有
\[
  \frac{f_y}{f_x} = \frac{n_x}{n_y}.
\]
因此，在已知 $f_x$ 的情况下，通过调节或判读稳定图形对应的 $n_x:n_y$，即可求出待测信号频率
\[
  f_y = \frac{n_x}{n_y} f_x.
\]

\subsection{数字示波器的数学运算}

数字示波器采集到两路信号后，可在内部对采样数据进行数学运算，例如通道相加、相减、相乘等。设两通道信号分别为 $u_1(t)$ 和 $u_2(t)$，加法运算显示
\[
  u_+(t)=u_1(t)+u_2(t),
\]
乘法运算显示
\[
  u_\times(t)=u_1(t)u_2(t).
\]
通过观察运算波形，可以直观理解双通道信号叠加、调制或乘积变化的特点。

\section{实验内容和步骤}

\subsection{观测函数发生器输出的波形}

\begin{enumerate}[label=\arabic*)]
  \item 将函数发生器 CH1 端输出信号接入数字示波器 CH1 通道，检查探头、接地端和通道设置。
  \item 分别调出正弦波、三角波和方波，使示波器稳定显示波形，并记录波形图与波形名称。
  \item 将函数发生器设置为 $100\,\si{\hertz}$ 的正弦波，使用示波器光标功能测量峰峰值、周期和频率，并记录测量图与结果。
  \item 对同一个 $100\,\si{\hertz}$ 正弦波使用菜单自动测量功能，记录峰峰值、周期和频率。
  \item 比较光标测量与自动测量结果的差异，完成阶段再根据实测数据进行分析。
\end{enumerate}

\begin{LabRecordTable}
  \centering
  \caption{函数发生器输出波形记录}
  \label{table:waveforms}
  \renewcommand{\arraystretch}{1.25}
  \setlength{\tabcolsep}{4pt}
  \begin{tabularx}{\textwidth}{C{2.2cm} Y Y Y}
    \toprule
    项目 & 波形一 & 波形二 & 波形三 \\
    \midrule
    \imagehead{3.2cm}{波形图} & \imagecell{3.2cm} & \imagecell{3.2cm} & \imagecell{3.2cm} \\
    波形名称 &  &  &  \\
    \bottomrule
  \end{tabularx}
\end{LabRecordTable}

\begin{LabRecordTable}
  \centering
  \caption{$100\,\si{\hertz}$ 正弦波的光标测量记录}
  \label{table:cursor-measure}
  \renewcommand{\arraystretch}{1.25}
  \setlength{\tabcolsep}{4pt}
  \begin{tabularx}{\textwidth}{C{2.4cm} Y}
    \toprule
    项目 & 记录内容 \\
    \midrule
    \imagehead{4.2cm}{测量图} & \imagecell{4.2cm} \\
    测量结果 & 峰峰值：\hspace{3cm} 周期：\hspace{3cm} 频率： \\
    \bottomrule
  \end{tabularx}
\end{LabRecordTable}

\begin{LabRecordTable}
  \centering
  \caption{$100\,\si{\hertz}$ 正弦波的自动测量记录}
  \label{table:auto-measure}
  \renewcommand{\arraystretch}{1.25}
  \setlength{\tabcolsep}{4pt}
  \begin{tabularx}{\textwidth}{C{2.4cm} Y}
    \toprule
    项目 & 记录内容 \\
    \midrule
    \imagehead{4.2cm}{测量图} & \imagecell{4.2cm} \\
    测量结果 & 峰峰值：\hspace{3cm} 周期：\hspace{3cm} 频率： \\
    \bottomrule
  \end{tabularx}
\end{LabRecordTable}

\subsection{利用李萨如图形测频率}

\begin{enumerate}[label=\arabic*)]
  \item 将函数发生器一路输出接入示波器 CH1 通道，作为 $X$ 方向信号；将另一路输出接入示波器 CH2 通道，作为 $Y$ 方向待测信号。
  \item 设置待测信号为一定频率的正弦信号，调节已知信号频率 $f_x$，使屏幕分别出现 $n_x:n_y=1:1$、$2:1$、$3:1$、$3:2$ 的稳定李萨如图形。
  \item 对每一种稳定图形拍照或记录，填写对应的 $f_x$、$f_y$，并由 $f_y=(n_x/n_y)f_x$ 计算待测频率。
\end{enumerate}

\begin{LabRecordTable}
  \centering
  \caption{李萨如图形测频率记录}
  \label{table:lissajous}
  \renewcommand{\arraystretch}{1.25}
  \setlength{\tabcolsep}{3pt}
  \begin{tabularx}{\textwidth}{C{2cm} *{4}{Y}}
    \toprule
    $n_x:n_y$ & $1:1$ & $2:1$ & $3:1$ & $3:2$ \\
    \midrule
    \imagehead{3cm}{图形} & \imagecell{3cm} & \imagecell{3cm} & \imagecell{3cm} & \imagecell{3cm} \\
    $f_x/\si{\hertz}$ &  &  &  &  \\
    $f_y/\si{\hertz}$ &  &  &  &  \\
    \bottomrule
  \end{tabularx}
\end{LabRecordTable}

\subsection{数字示波器的数学运算}

\begin{enumerate}[label=\arabic*)]
  \item 设置函数发生器两路输出信号，将两路信号分别接入示波器 CH1 和 CH2 通道。
  \item 打开示波器数学运算功能，选择 CH1+CH2，调节显示比例并记录运算波形。
  \item 改变两路信号参数，选择 CH1*CH2，调节显示比例并记录运算波形。
  \item 比较两种数学运算波形与原始双通道波形的关系。
\end{enumerate}

\begin{LabRecordTable}
  \centering
  \caption{数字示波器数学运算记录}
  \label{table:math-operation}
  \renewcommand{\arraystretch}{1.25}
  \setlength{\tabcolsep}{4pt}
  \begin{tabularx}{\textwidth}{C{2.4cm} Y Y}
    \toprule
    项目 & CH1+CH2 & CH1*CH2 \\
    \midrule
    \imagehead{4cm}{运算图} & \imagecell{4cm} & \imagecell{4cm} \\
    信号设置 &  &  \\
    波形特征 &  &  \\
    \bottomrule
  \end{tabularx}
\end{LabRecordTable}

\section{实验数据处理}

本报告处于准备阶段，待完成实验并确认原始数据后，再根据表 \ref{table:cursor-measure}、表 \ref{table:auto-measure}、表 \ref{table:lissajous} 和表 \ref{table:math-operation} 中的数据补充计算过程、结果比较和误差分析。

\section{实验的分析和总结}

本报告处于准备阶段，待完成实验并确认原始数据后，再结合实际波形、测量结果和李萨如图形记录完成分析和总结。

\insertRawDataPage

\end{document}
